\input{../tex-inputs/latex-leseansicht-vorspann}

\section[Arthur Schnitzler an Gustav Schwarzkopf, {[}zwischen 1899 und 1919?{]}]{L04177 Arthur Schnitzler an Gustav Schwarzkopf, {[}zwischen 1899 und 1919?{]}}
\nopagebreak\mylabel{L04177v}
\rehead{ }\normalsize\beginnumbering\briefempfaengerindex{Schwarzkopf, Gustav@\textsc{Schwarzkopf, Gustav}!zzzSchnitzler, Arthur@\emph{von Arthur Schnitzler}!1919-12-311@{{[}zwischen März 1899 und 1931?{]}}|(be}
\toendnotes[C]{\smallbreak\pagebreak[2]}
\correspDesc{Versand  durch Arthur Schnitzler im Zeitraum [zwischen März 1899 und 1931?] in Wien
\newline{}Erhalt  durch Gustav Schwarzkopf im Zeitraum [zwischen März 1899 und 1931?] in Wien}\toendnotes[C]{\smallbreak}
\Standort{CUL, Schnitzler, B 96.}
\physDesc{Brief, 1 Blatt, 2 Seiten, 179 Zeichen
\newline{}Handschrift: Bleistift, deutsche Kurrent}\toendnotes[C]{\smallbreak}
\pstart
           \noindent{}{\pb}\label{K_L04177-1v}\edtext{lieber Guſtav}{\lemma{\textnormal{\emph{lieber Gustav}}}\Cendnote{\textnormal{Das Korrespondenzstück ist undatiert. Die Anrede »lieber Guſtav« verwendete er aber est ab
                     XXXX Auszeichnungsfehler: Dokument L04122 nicht gefunden, so dass zeitlich eine Grenze 
                     nach vorne möglich ist. Nach hinten dürfte die Praxis des Theaterkarten-Zusendens spätestens um 1919 ein Ende gefunden zu haben,
                     vgl. XXXX Auszeichnungsfehler: Dokument L04159 nicht gefunden.}}}\label{K_L04177-1}, ich
               hoffe Sie ſind in der Lage von dieſem Logenanweiſg Gebrauch zu machen, vielleicht mit
               Ihrem Bruder\pwindex{Schwarzkopf, Max 12.\,6.\,1857 Wien – 14.\,4.\,1928 ebd.@\textsc{Schwarzkopf, Max} (12.\,6.\,1857 Wien – 14.\,4.\,1928 ebd.)|pwv}? Zurückſendg an
               mich in keinem Falle {\pb}nöthig.\pend
           
\pstart
           Herzlichſt Ihr{\\[\baselineskip]}\spacefill\mbox{Arthur}\pend
           \leftskip=0em{}\selectlanguage{ngerman}\endnumbering\briefempfaengerindex{Schwarzkopf, Gustav@\textsc{Schwarzkopf, Gustav}!zzzSchnitzler, Arthur@\emph{von Arthur Schnitzler}!1899-03-011@{{[}zwischen März 1899 und 1931?{]}}|)be}\mylabel{L04177h}
\begin{anhang}
\end{anhang}\newcommand{\dateiname}{L04177}\newcommand{\titel}{Arthur Schnitzler an Gustav Schwarzkopf, [zwischen 1899 und 1919?]}\newcommand{\editorInnen}{Herausgegeben von Jahnke, SelmaMüller, Martin Anton}\input{../tex-inputs/latex-leseansicht-abspann}
